% !TeX root = ../0_Manuscript.tex

\section{BBI in practice \ddcpu}
\label{chap:2_goodPractices;sect:bbiInPractice}
With actual \bbi platform in mind, allow me to introduce a typical \bbi platform model that we are going to study in details.
%In this section, I am going to describe coarse electrical \bbi platform models.
These models allow us to concisely understand, evaluate and simulate a \bbi platform's behavior, limitations, and room for improvement.
We will therefore analyze the platform's performance and point out its weaknesses.

    \subsection{Typical BBI platform model \ddcu}
    \label{chap:2_goodPractices;sect:bbiInPractice;subsect:bbiPlatformModel}
    % !TeX root = ../0_Manuscript.tex

%\begin{figure}[H]
%    \centering
%    \includegraphics[width=0.7\textwidth, center]{2_goodPractices/figures/model0Probe.pdf}
%    \caption{Simplified model of a \bbi platform in the state-of-the-art. The model represents the major tools involved in a \bbi platform, with the voltage pulse generator, the \bbi metal probe, the transmission line used to interconnect the generator to the probe, the IC, modeled with a 1 k\textOmega\xspace resistor, and the platform grounding impedance, represented with two 150 \textOmega\xspace resistors. This model allows for fast evaluation and prediction of the \bbi platform macro-behavior.}
%    \label{fig:bbiPracticeBadGnd}
%\end{figure}

\begin{figure}[ht]
    \centering
    \begin{subfigure}{0.45\textwidth}
        \includegraphics[width=1.0\textwidth, center]{2_goodPractices/figures/model0Probe.pdf}
        \caption{Model of a state-of-the-art \bbi setup, with its key components.}
        \label{fig:bbiPracticeBadGnd}
    \end{subfigure}
    \hspace{0.045\textwidth}
    \begin{subfigure}{0.45\textwidth}
        \includegraphics[width=1.0\textwidth, center]{2_goodPractices/figures/SIMPLE_BBI_0.pdf}
        \caption{Simulation results of the state-of-the-art \bbi platform model.}
        \label{fig:bbiPracticeBadGndSignals}
    \end{subfigure}
    \caption{State-of-the-art \bbi platform electrical model and its simulation results.}
    \label{fig:bbiBadGndGlobalFig}
\end{figure}

    To be able to quickly predict and analyze \bbi platforms, I developed an elementary electrical model, illustrated in Fig. \ref{fig:bbiPracticeBadGnd}.
    This model represents the key components of a \bbi platform, which are:
    \begin{itemize}
        \setlength\itemsep{-0.1em}
        \item The voltage pulse generator;
        \item The transmission line, used to connect the probe to the generator;
        \item The \bbi probe;
        \item The targeted IC, modeled by an electrical resistance;
        \item The grounding installation, consisting of electrical resistances connected between equipment grounding.
    \end{itemize}


    \subsection{Platforms evaluation criteria \ddcu}
    \label{chap:2_goodPractices;sect:bbiInPractice;subsect:platformEvalCrit}
    For the purpose of evaluating \bbi platforms, we decided to focus on two important criteria, allowing to represent the platform quality:
    \begin{itemize}
        \setlength\itemsep{-0.1em}
        \item The characteristics of the voltage pulse measured at the generator output, allowing to observe how the generator behaves when loaded with the transmission line and the IC;
        \item The characteristics of the target ground current waveform, allowing to monitor exactly what is actually injected into the IC.
    \end{itemize}

%        \subsubsection{Voltage pulse generator output characteristics measurements \ddc}
%        \label{chap:2_goodPractices;sect:bbiInPractice;subsect:platformEvalCrit;subsubsect:vOutputChar}
%
%        \subsubsection{IC target ground current measurements \ddc}
%        \label{chap:2_goodPractices;sect:bbiInPractice;subsect:platformEvalCrit;subsubsect:icGndMeas}

    \subsection{Raw results \ddcu}
    \label{chap:2_goodPractices;sect:bbiInPractice;subsect:rawRes}
    To that end, we will deeply analyze the platform's simulation results shown in Fig. \ref{fig:bbiPracticeBadGndSignals}.
    Four signals are displayed, with their colors matching the colors in Fig. \ref{fig:bbiPracticeBadGnd} for greater clarity.
    There are three voltage waveforms and a current waveform.
    The blue waveform is the ideal voltage pulse an attacker wants to apply to the backside of an IC during a body biasing injection.
    Its characteristics are the following: a voltage set point of -140 V and a pulse width set point of 10 ns.
    It is a steep, fast and precise pulse with controlled rise and fall times, pulse width and voltage.
    However, when performing real experiments, which the model allows us to evaluate, this ideal pulse falls apart.
    It can be seen thanks to the orange and green waveforms, representing respectively the pulse observer at the generator output and the pulse effectively applied onto the backside of the IC target.
    There are multiple obvious observations that can be made concerning the received pulse (green) signal:
    \begin{itemize}
        \setlength\itemsep{-0.1em}
        \item The voltage set point is not respected, with a 23.5 \% negative percentage overshoot (PO) on the falling edge, and a 107 \% positive percentage overshoot (PO);
        \item There is obvious ringing, causing the pulse width to be longer than expected in addition to damped oscillations
    \end{itemize}
    These effects can also be noticed on the IC ground current waveform (\textcolor{Purple}{purple}), as it is a mirror of the applied pulse due to the pure resistive nature of the IC in that model.

    \subsection{Analysis conclusions \ddcu}
    \label{chap:2_goodPractices;sect:bbiInPractice;subsect:analysConcl}
    To comprehend the implications of such observations, it is imperative that we thoroughly analyze each one of them.

    The first important thing to note is that the various model numeric values are extracted from our actual platform.
    Therefore, the 150 \textOmega\xspace grounding and the 1 k\textOmega\xspace IC are average measured values of actual devices.
    Thus, these parameters, in addition to the transmission line characteristics, will inevitably vary with a certain amount from one platform to another.
    
\begin{figure}[ht]
    \centering
    \begin{subfigure}{0.48\textwidth}
        \centering
        \includegraphics[width=1.0\textwidth, center]{2_goodPractices/figures/SIMPLE_BBI_0_0_IC250.pdf}
        \caption{Simulation result with an IC load equals to 250 \textOmega\xspace. The voltage set point is not met, and the amplitude is -30 \% lower than requested, with obvious ringing. Due to the lower IC load value, there is more energy injected into the IC. This is mainly caused by the voltage divider formed by the generator output stage and the IC load.}
        \label{fig:bbiPracticeBadGndICLoad0}
    \end{subfigure}
    \hfill
    \begin{subfigure}{0.48\textwidth}
        \centering
        \includegraphics[width=1.0\textwidth, center]{2_goodPractices/figures/SIMPLE_BBI_0_0_IC2k.pdf}
        \caption{Simulation result with an IC load equals to 2 k\textOmega\xspace. The voltage set point is not met once again, and the amplitude is 40 \% higher than the requested value, with obvious ringing on all waveforms. Due to the higher IC load value, there is less energy injected into the IC. This is mainly caused by the voltage divider formed by the generator output stage and the IC load.}
        \label{fig:bbiPracticeBadGndICLoad1}
    \end{subfigure}
%    \vskip\baselineskip◘
%    \par\bigskip
%    \begin{subfigure}{0.48\textwidth}
%        \centering
%        \includegraphics[width=1.0\textwidth, center]{2_goodPractices/figures/SIMPLE_BBI_0_0_IC1k.pdf}
%        \caption{Caption here.}
%        \label{fig:bbiPracticeBadGndICLoad2}
%    \end{subfigure}
%    \hfill
%    \begin{subfigure}{0.48\textwidth}
%        \centering
%        \includegraphics[width=1.0\textwidth, center]{2_goodPractices/figures/SIMPLE_BBI_0_0_IC2k.pdf}
%        \caption{Caption here.}
%        \label{fig:bbiPracticeBadGndICLoad3}
%    \end{subfigure}
    \caption{Platform simulation results with different IC load values.}
    \label{fig:bbiBadGndIcLoadVar}
\end{figure}

    Indeed, the backside surface of an IC does not equal to a constant load.
    In addition to this, if the IC substrate is thinned, these values will change even more.
    Therefore, depending on the probe location and the IC substrate thickness, the generator might not see the same load.
    To illustrate the induced effects of such differences, I performed simulations with various IC values, representing typical measured values for my IC target when thinned down to 50 µm, up to more than 700 µm, and the results are shown in Fig. \ref{fig:bbiBadGndIcLoadVar}, both for a 250 \textOmega\xspace load (Fig. \ref{fig:bbiPracticeBadGndICLoad0}), and a 2 k\textOmega\xspace load (Fig. \ref{fig:bbiPracticeBadGndICLoad1}).
    In both cases, due to the non-zero generator output impedance, the latter forms a voltage divider with the IC load.
    On the one hand, with an IC load value one quarter lower, there is more current in it, while the applied pulse amplitude is 30\% lower.
    On the other hand, with an IC load value two times higher, there is less current in it, while the applied pulse amplitude is 40\% higher.
    It represents a 70 \% range around the set point value, which is excessively high.
    However, in both cases, the ringing is still present with the same amount relative to the pulse amplitude.

    Eventually, all of these observations allow us to spot three major flaws of such a platform:
    \begin{itemize}
        \setlength\itemsep{-0.1em}
        \item The platform parameters are difficult to control, leading to unknown values concerning pulse width, voltage set point, etc.;
        \item It leads to a poor temporal accuracy, thus minimizing the chances to perform a precise and repeatable fault injection;
        \item At last, all parameters are platform-dependent, leading to a low reproducibility rate, thus lowering the credibility of experiments performed on such platforms.
    \end{itemize}
    In this context, I present in the next section various simple improvements to the \bbi state-of-the-art platform.
